\section{The Width-3 Gadget}

\begin{frame}{The Width-3 Gadget}

\begin{definition}
The previous graph gives rise to the following width-3 gadget:
\begin{align*}
\begin{quantikz}
\lstick{$q_1$} 
    & \gate{R_z(\theta_1)} & \targ{}   & \ctrl{1} & \gate{H} & \qw \\
\lstick{$q_2$} 
    & \gate{R_z(\theta_2)} & \qw       & \ctrl{-1}& \gate{H} & \qw \\
\lstick{$q_3$} 
    & \gate{R_z(\theta_3)} & \ctrl{-2} & \qw      & \gate{H} & \qw
\end{quantikz}        
\end{align*}
\end{definition}

Question: Is this gadget universal for MBQC?

\end{frame}

% \begin{frame}{Linear Algebra Universality Algorithm}

% \begin{theorem}
% A gate set $\mathcal{G}$ is universal for $n$ qubits if it generates $SU(2^n)$, the group of $2^n \times 2^n$ unitary matrices with determinant 1.
% \end{theorem}

% \begin{theorem}[Sawicki and Karnas, 2017]
% Given a finite gate set $\mathcal{G}$, we can determine if it is universal in finite time by computing the kernel, trace, and product of some matrices and vectors.
% \end{theorem}

% \begin{problem}<2->
% The gate set $\mathcal{G}$ from the width-3 gadget is not finite, since we can choose $\theta_i \in [0, 2\pi]$.
% \end{problem}

% \end{frame}

\begin{frame}{Lie Algebra Universality Algorithm}

\begin{definition}[Matrix Exponentiation]
For any $A \in \mathbb{C}^{m \times m}$, 
\begin{align*}
    e^A = \sum_{k=0}^\infty \frac{A^k}{k!}
\end{align*}
\end{definition}

\begin{definition}<2->[Lie Algebra of $SU(2^n)$]
\begin{align*}
    \mathfrak{su}(2^n) = \{ -iH \in \mathbb{C}^{2^n \times 2^n} \mid H^\dagger = H, \text{tr}(H) = 0 \}
\end{align*}
\end{definition}

\end{frame}

\begin{frame}{Lie Algebra Universality Algorithm cont'd}

\begin{theorem}
A gate set $\mathcal{G}$ is universal for $n$ qubits if it generates $SU(2^n)$, the group of $2^n \times 2^n$ unitary matrices with determinant 1.
\end{theorem}

\begin{fact}<2->
$U \in SU(2^n)$ if and only if $U = e^{H}$ for some $H \in \mathfrak{su}(2^n)$.
\end{fact}

\begin{corollary}<2->
To show $\mathcal{G}$ is universal, it suffices to generate the Lie algebra.
\end{corollary}

\begin{fact}<3->
$H \in \mathfrak{su}(2^n)$ if and only if $H = -i \sum_{P \in \mathcal{P}_n \setminus \{I\}} \alpha_P P$ for some $\alpha_P \in \mathbb{R}$, where $\mathcal{P}_n$ is the set of Pauli matrices on $n$ qubits.
\end{fact}

\begin{corollary}<3->
The traceless Hermitian operators are exactly those that can be represented as a linear combination of Pauli operators (X, Y, Z).
\end{corollary}

\end{frame}

\begin{frame}{Lie Algebra Universality Algorithm: Cliffords}

\begin{lemma}
Denote by $G(\theta_1, \theta_2, \theta_3)$ the width-3 gadget. Then $G(0,0,0)$ is some fixed Clifford operation $C$ and $C \in \mathcal{G}$. 
\end{lemma}

\begin{lemma}<2->
If $C \in \mathcal{G}$ then $C^\dagger \in \mathcal{G}$.
\end{lemma}

\begin{proof}<3->
The Clifford group on 3 qubits is finite. Hence every element has finite order, i.e. $\exists t$ with $C^t = I$. Thus $C^\dagger = C^{t-1} \in \mathcal{G}$.
\end{proof}

\end{frame}

\begin{frame}{Lie Algebra Universality Algorithm: Addition}

\begin{fact}[Baker-Campbell-Hausdorff formula]
For $A,B \in \mathfrak{su}(2^n)$,
\begin{align*}
    e^{A} e^{B} = e^{A + B + \frac{1}{2}[A,B] + \frac{1}{12}[A,[A,B]] + \frac{1}{12}[B,[A,B]] + \cdots}
\end{align*}
where $[A,B] = AB - BA$ is the commutator of $A$ and $B$.
\end{fact}

\begin{fact}
For $A,B \in \mathfrak{su}(2^n)$,
\begin{align*}
    \lim_{k \to \infty} e^{A/k} e^{B/k} = e^{A + B}
\end{align*}
\end{fact}

\begin{corollary}
If we have $A,B$ in the Lie algebra, we can generate their sum.
\end{corollary}

\end{frame}

\begin{frame}{Lie Algebra Universality Algorithm: Commutators}

\begin{fact}
For $A,B \in \mathfrak{su}(2^n)$,
\begin{align*}
    \lim_{k \to \infty} e^{B/k} e^{A/k} e^{-B/k} e^{-A/k} = e^{[A,B]}.
\end{align*}
\end{fact}

\begin{corollary}
If we have $A,B$ in the Lie algebra, we can generate their commutator $[A,B]$.
\end{corollary}

\end{frame}

\begin{frame}{Lie Algebra Universality Algorithm: Operations}

Strategy: Start with some elements in the Lie algebra and generate more using:
\begin{itemize}
    \item Conjugation by $C$
    \item Addition
    \item Commutators
\end{itemize}

\begin{fact}<2->
We cannot generate all $|\{I,X,Y,Z\}|^3 - |(I,I,I)| = 63$ non-identity Pauli operators on $3$ qubits this way.
\end{fact}

\begin{proof}<3->
Brute force using a computer program.
\end{proof}

\end{frame}

\begin{frame}{Lie Algebra Universality Algorithm: Visual Intuition}

The $CNOT$ ``forbids us'' from generating $R_x(\theta_1)$ alone on $q_1$.
\begin{center}
\begin{tikzpicture}[scale=0.7, transform shape]

\node (left) at (0,0) {
\begin{quantikz}
\lstick{$q_1$} 
    & \gate{R_z(\theta_1)} & \targ{}   & \ctrl{1} & \gate{H} & \qw \\
\lstick{$q_2$} 
    & \gate{R_z(\theta_2)} & \qw       & \ctrl{-1}& \gate{H} & \qw \\
\lstick{$q_3$} 
    & \gate{R_z(\theta_3)} & \ctrl{-2} & \qw      & \gate{H} & \qw
\end{quantikz}
};

\node (eq) at (3.25,0) {$=$};

\node (right) at (7,0) {
\begin{quantikz}
\lstick{$q_1$} 
    & \targ{}   & \ctrl{1} & \gate{H} & \gate{R_x(\theta_1)R_x(\theta_3)} & \qw \\
\lstick{$q_2$} 
    & \qw       & \ctrl{-1}& \gate{H} & \gate{R_x(\theta_2)} & \qw \\
\lstick{$q_3$} 
    & \ctrl{-2} & \qw      & \gate{H} & \gate{R_x(\theta_3)} & \qw
\end{quantikz}
};

\end{tikzpicture}
\end{center}

\only<2->{
Idea: ``Stitch'' this with another gadget to ``get'' $R_x(\theta_1)$ on $q_1$.
}

\end{frame}